% GENERATED FILE. Edit canonical lessons, then run npm run generate:books.
% canonical-source-hash: fnv1a64:18b32d35be10cdeb
% canonical-lessons: MR-C112-dur, MR-C112-at, MR-C112-baher, MR-C112-mage, MR-C112-pudhe

\chapter{Far, Inside, Outside, Behind, Ahead}
\label{ch:mr-far-inside-outside-behind-ahead}
\hlchaptermodality{112}

\begin{chapteropening}
\textbf{By the end of this chapter:} \emph{I can say far, inside, outside, behind and ahead, in front.}

Complete the last lesson of chapter 112: I can say far, inside, outside, behind and ahead, in front.
\end{chapteropening}

\section[\texorpdfstring{\mr{दूर} (dūr)}{दूर (dūr)}]{\mr{दूर} (dūr) — far}

\label{lesson:MR-C112-dur}



\begin{quote}
Before the new one: say the Marathi for after, then the Marathi for time.
\end{quote}

\subsection*{You'll want to know: \mr{दूर}}
\textbf{\mr{दूर}} — \emph{dūr} — \textquotedblleft{}far\textquotedblright{}.

\begin{cousinweb}
\textbf{\mr{दूर}} (\emph{dūr}) is Sanskrit \emph{dūra}, \textquotedblleft{}distant\textquotedblright{}, which has lost only its final vowel (CDIAL 6495).
\end{cousinweb}

\begin{grammarlens}[title={What you've built}]
One more word for this chapter.
\end{grammarlens}

\subsection*{Your turn}
\begin{itemize}
\raggedright
  \item \emph{Say it:} \emph{dūr}
  \item \emph{Say it:} \emph{dūr}, once more
  \item \emph{Say it:} say \emph{dūr}
  \item \emph{Recall:} say the Marathi for after, then the Marathi for time, then say \emph{dūr} again
\end{itemize}

\begin{tcolorbox}[breakable,colback=teal!4,colframe=teal!35!black,title={Before you move on}]
What does \mr{दूर} mean? (\textquotedblleft{}Far\textquotedblright{}.) Say it once more.
\end{tcolorbox}
\section[\texorpdfstring{\mr{आत} (āt)}{आत (āt)}]{\mr{आत} (āt) — inside}

\label{lesson:MR-C112-at}



\begin{quote}
Before the new one: say the Marathi for time, then the Marathi for far.
\end{quote}

\subsection*{You'll want to know: \mr{आत}}
\textbf{\mr{आत}} — \emph{āt} — \textquotedblleft{}inside\textquotedblright{}.

\begin{cousinweb}
\textbf{\mr{आत}} (\emph{āt}) is a tadbhava, a word inherited from Sanskrit and reshaped by sound change. It comes from \emph{antar}, \textquotedblleft{}inside\textquotedblright{} (CDIAL 356).
\end{cousinweb}

\begin{grammarlens}[title={What you've built}]
One more word for this chapter.
\end{grammarlens}

\subsection*{Your turn}
\begin{itemize}
\raggedright
  \item \emph{Say it:} \emph{āt}
  \item \emph{Say it:} \emph{āt}, once more
  \item \emph{Say it:} say \emph{āt}
  \item \emph{Recall:} say the Marathi for time, then the Marathi for far, then say \emph{āt} again
\end{itemize}

\begin{tcolorbox}[breakable,colback=teal!4,colframe=teal!35!black,title={Before you move on}]
What does \mr{आत} mean? (\textquotedblleft{}Inside\textquotedblright{}.) Say it once more.
\end{tcolorbox}
\section[\texorpdfstring{\mr{बाहेर} (bāher)}{बाहेर (bāher)}]{\mr{बाहेर} (bāher) — outside}

\label{lesson:MR-C112-baher}



\begin{quote}
Before the new one: say the Marathi for far, then the Marathi for inside.
\end{quote}

\subsection*{You'll want to know: \mr{बाहेर}}
\textbf{\mr{बाहेर}} — \emph{bāher} — \textquotedblleft{}outside\textquotedblright{}.

\begin{cousinweb}
\textbf{\mr{बाहेर}} (\emph{bāher}) is a tadbhava, a word inherited and reshaped by sound change. It comes from \emph{*bāhira}, \textquotedblleft{}external\textquotedblright{}, a form CDIAL reconstructs from its descendants (CDIAL 9226). CDIAL builds it on Sanskrit \emph{bahis}, \textquotedblleft{}outside\textquotedblright{}.
\end{cousinweb}

\begin{grammarlens}[title={What you've built}]
One more word for this chapter.
\end{grammarlens}

\subsection*{Your turn}
\begin{itemize}
\raggedright
  \item \emph{Say it:} \emph{bāher}
  \item \emph{Say it:} \emph{bāher}, once more
  \item \emph{Say it:} say \emph{bāher}
  \item \emph{Recall:} say the Marathi for far, then the Marathi for inside, then say \emph{bāher} again
\end{itemize}

\begin{tcolorbox}[breakable,colback=teal!4,colframe=teal!35!black,title={Before you move on}]
What does \mr{बाहेर} mean? (\textquotedblleft{}Outside\textquotedblright{}.) Say it once more.
\end{tcolorbox}
\section[\texorpdfstring{\mr{मागे} (māge)}{मागे (māge)}]{\mr{मागे} (māge) — behind}

\label{lesson:MR-C112-mage}



\begin{quote}
Before the new one: say the Marathi for inside, then the Marathi for outside.
\end{quote}

\subsection*{You'll want to know: \mr{मागे}}
\textbf{\mr{मागे}} — \emph{māge} — \textquotedblleft{}behind\textquotedblright{}.

\begin{cousinweb}
The origin of \textbf{\mr{मागे}} (\emph{māge}) is not settled. Neither CDIAL nor DEDR has an entry for it. Marathi \emph{māg}, \textquotedblleft{}road, track\textquotedblright{}, which CDIAL lists under Sanskrit \emph{mārga}, \textquotedblleft{}track\textquotedblright{} (CDIAL 10071), is a possible lead, but no source links the two.
\end{cousinweb}

\begin{grammarlens}[title={What you've built}]
One more word for this chapter.
\end{grammarlens}

\subsection*{Your turn}
\begin{itemize}
\raggedright
  \item \emph{Say it:} \emph{māge}
  \item \emph{Say it:} \emph{māge}, once more
  \item \emph{Say it:} say \emph{māge}
  \item \emph{Recall:} say the Marathi for inside, then the Marathi for outside, then say \emph{māge} again
\end{itemize}

\begin{tcolorbox}[breakable,colback=teal!4,colframe=teal!35!black,title={Before you move on}]
What does \mr{मागे} mean? (\textquotedblleft{}Behind\textquotedblright{}.) Say it once more.
\end{tcolorbox}
\section[\texorpdfstring{\mr{पुढे} (puḍhe)}{पुढे (puḍhe)}]{\mr{पुढे} (puḍhe) — ahead, in front}

\label{lesson:MR-C112-pudhe}



\begin{quote}
Before the new one: say the Marathi for outside, then the Marathi for behind.
\end{quote}

\subsection*{You'll want to know: \mr{पुढे}}
\textbf{\mr{पुढे}} — \emph{puḍhe} — \textquotedblleft{}ahead, in front\textquotedblright{}.

\begin{cousinweb}
The origin of \textbf{\mr{पुढे}} (\emph{puḍhe}) is not settled. Neither CDIAL nor DEDR, searched by form and by meaning, has an entry for it.
\end{cousinweb}

\begin{grammarlens}[title={What you've built}]
One more word for this chapter.
\end{grammarlens}

\subsection*{Your turn}
\begin{itemize}
\raggedright
  \item \emph{Say it:} \emph{puḍhe}
  \item \emph{Say it:} \emph{puḍhe}, once more
  \item \emph{Say it:} say \emph{puḍhe}
  \item \emph{Recall:} say the Marathi for outside, then the Marathi for behind, then say \emph{puḍhe} again
\end{itemize}

\begin{tcolorbox}[breakable,colback=teal!4,colframe=teal!35!black,title={Before you move on}]
What does \mr{पुढे} mean? (\textquotedblleft{}Ahead, in front\textquotedblright{}.) Say it once more.
\end{tcolorbox}
